\section{Problem Statement}
\label{sec:problem}

Receipt key-information extraction (KIE) takes a scanned receipt
image $I$ and returns a structured record
$\hat{y} = (\hat{y}_{\textsf{company}}, \hat{y}_{\textsf{date}},
\hat{y}_{\textsf{address}}, \hat{y}_{\textsf{total}})$ of four
free-text field values.  We evaluate against the canonical
ICDAR-2019 SROIE Task-3 test split~\cite{huang2019icdar} of $n=347$
Malaysian retail receipts.

The two classes of system we compare differ in how the field
assignment is made.  An end-to-end model
(DONUT~\cite{kim2022donut}) decodes the structured output
autoregressively from the image; a modular pipeline
($\textsf{detect} \to \textsf{read} \to \textsf{assign}$) detects
text-line bounding boxes (YOLOv8~\cite{jocher2023yolov8}), reads each
line (TrOCR~\cite{li2023trocr}), and assigns lines to fields by
deterministic regular expressions plus a small generative segmentation
prior over the receipt's vertical structure.

The metric is whitespace-token F1 averaged uniformly over the four
fields ($F_1^{\text{macro}}$), reported alongside Normalised Edit
Distance (NED) and strict Exact Match (EM).  Statistical uncertainty
is quantified by paired-bootstrap 95\,\% confidence intervals over
per-image correctness vectors and McNemar's exact test on per-image
all-fields-correct outcomes.  All inputs are normalised through the
symmetric per-field map of \texttt{models/normalize\_bundle.py} so the
two arms see the same canonical form for each prediction and ground-
truth pair.
